\section{幂级数}
\label{sec:ru-power-series}

在一般的函数项级数中，有一类各项由同一中心的整数次幂组成：
\[
 \sum_{k=0}^\infty c_k(x-x_0)^k,
\]
称为以$x_0$为中心的\term{幂级数}（power series），$c_k$称为系数。它的部分和是多项式，多项式只用加法和乘法就能计算，这使幂级数成为表示和近似函数的基本工具。
本节研究幂级数在哪些输入处收敛，以及收敛怎样控制和函数的连续性。

前面已经遇到两个幂级数：几何级数$\sum_kx^k$在$|x|<1$时收敛到$1/(1-x)$，指数级数$\sum_kx^k/k!$在每个实数处收敛。对收敛的输入，级数的和定义了一个函数，称为和函数。
几何级数在$|x|\geq1$时发散，可见幂级数未必处处收敛。下面先确定收敛的输入范围，再研究在什么范围内能对全部输入统一控制误差。已有函数怎样展开成幂级数，留到第~\ref{sec:ru-elementary-taylor}节。

\subsection{收敛半径与收敛区间}
\label{sec:ru-radius-of-convergence}

先看一点处的收敛对其他输入意味着什么。下面的结论说明，幂级数的收敛范围总是关于中心对称的一段（可能还包括端点）。
\begin{theorem}[Abel第一定理]
\label{thm:ru-abel-lemma}
若幂级数$\sum c_k(x-x_0)^k$在某点$x_1\neq x_0$处收敛，记$r=|x_1-x_0|$，则它在每个满足$|x-x_0|<r$的$x$处绝对收敛。
\end{theorem}
\begin{proof}
$x_1$处的级数收敛，所以其项趋于零，进而有界：存在$M$使$|c_k|r^k\leq M$对所有$k$成立。对$|x-x_0|<r$，记$q=|x-x_0|/r<1$，则
\[
 |c_k(x-x_0)^k|=|c_k|r^kq^k\leq Mq^k,
\]
与收敛的几何级数$\sum Mq^k$比较，即得绝对收敛。
\end{proof}
取逆否命题：若级数在某点$x_2$处发散，则它在每个离中心更远的点$x$（$|x-x_0|>|x_2-x_0|$）处也发散；否则那一点的收敛会推出$x_2$处收敛。

于是令$R=\sup\{|x-x_0|:\text{级数在}x\text{处收敛}\}\in[0,+\infty]$，称为\term{收敛半径}（radius of convergence）。由Abel第一定理，$|x-x_0|<R$时级数绝对收敛，$|x-x_0|>R$时级数发散。
$R=0$时级数只在中心处收敛；$R=+\infty$时在整条实轴上收敛；$0<R<+\infty$时，开区间$(x_0-R,x_0+R)$称为收敛区间，在其中绝对收敛，在其外发散。
两个端点$x_0\pm R$处的行为定理没有给出，必须各自检验。实际收敛的输入组成的集合称为收敛域，它是收敛区间加上收敛的端点。

\paragraph{收敛半径的计算}
收敛半径可以由系数算出。对固定的$x$，把根值判别用于$\sum|c_k|\,|x-x_0|^k$：$\limsup(|c_k|\,|x-x_0|^k)^{1/k}=|x-x_0|\limsup|c_k|^{1/k}$，它小于$1$时收敛，大于$1$时发散。由此得到
\[
 \frac1R=\limsup_{k\to\infty}|c_k|^{1/k},
\]
称为Cauchy--Hadamard公式，约定右侧为$0$时$R=+\infty$、为$+\infty$时$R=0$。
若系数非零且$|c_{k+1}/c_k|\to\ell$，用比值判别同样得到$R=1/\ell$，计算往往更方便。比值算法要求相邻系数非零且比值有极限；不满足时仍可用Cauchy--Hadamard公式。
例如$\sum x^{2k}$的奇数次系数为零，相邻系数之比没有定义，但$|c_k|^{1/k}$在偶数$k$处为$1$、奇数$k$处为$0$，上极限为$1$，所以$R=1$。

几个例子说明公式怎样用，也说明半径不能代替端点判别。指数级数的系数比为$1/(k+1)\to0$，$R=+\infty$；$\sum k!\,x^k$的系数比$k+1\to\infty$，$R=0$，只在$x=0$处收敛。
$\sum x^k$、$\sum x^k/k$、$\sum x^k/k^2$的收敛半径都是$1$，端点行为却各不相同：第一个在$x=\pm1$处都发散，收敛域为$(-1,1)$；
第二个在$x=1$处是调和级数，发散，在$x=-1$处是交错调和级数的相反数，条件收敛，收敛域为$[-1,1)$；第三个在两端都绝对收敛，收敛域为$[-1,1]$。

\subsection{区间内的一致收敛性}
\label{sec:ru-power-series-uniform}

在整个收敛区间上未必能统一控制误差。几何级数截断到第$N$项的误差为$|x|^{N+1}/|1-x|$，在$[-\rho,\rho]$（$\rho<1$）上不超过$\rho^{N+1}/(1-\rho)$，可以用一个$N$同时控制；
在整个$(-1,1)$上，$x\to1^-$时误差趋于无穷，不存在共同的$N$。退一步，在严格位于收敛区间内部的闭子区间上，统一控制总是可以做到的。

\paragraph{内闭一致收敛}
设$R>0$，取$0<\rho<R$。在$|x-x_0|\leq\rho$上，$|c_k(x-x_0)^k|\leq|c_k|\rho^k$；而$x_0+\rho$位于收敛区间内部，级数在那里绝对收敛，即$\sum|c_k|\rho^k<\infty$。
由Weierstrass判别法，幂级数在$[x_0-\rho,x_0+\rho]$上一致收敛。由于$\rho$可以任意接近$R$，收敛区间内的每个闭子区间都含于某个这样的区间，这称为内闭一致收敛。
它不等于在整个开区间上一致收敛：几何级数在$(-1,1)$上的最大误差为无穷。

\paragraph{端点处的一致收敛}
端点的情形需要另外处理。若$0<R<\infty$且级数在端点$x_0+R$处收敛，那么在从中心到这个端点的整个闭区间上都可以统一控制误差，即使端点处只是条件收敛。
\begin{theorem}[Abel第二定理]
\label{thm:ru-abel-endpoint}
设$0<R<\infty$。若$\sum c_kR^k$收敛，则幂级数$\sum c_k(x-x_0)^k$在$[x_0,x_0+R]$上一致收敛。另一端$x_0-R$有对称的结论；两端都收敛时，在整个$[x_0-R,x_0+R]$上一致收敛。
\end{theorem}
\begin{proof}
令$t=(x-x_0)/R\in[0,1]$，$a_k=c_kR^k$，只需证明$\sum a_kt^k$在$[0,1]$上一致收敛。记$A_{n,m}=\sum_{k=n}^ma_k$。由$\sum a_k$收敛与柯西准则，给定$\varepsilon>0$，存在$N$使$m\geq n\geq N$时$|A_{n,m}|<\varepsilon$。
用第~\ref{sec:ru-series-tests}节的有限分部求和，令$A_{n,n-1}=0$，
\[
 \sum_{k=n}^ma_kt^k=A_{n,m}t^m+\sum_{k=n}^{m-1}A_{n,k}(t^k-t^{k+1}).
\]
$t\in[0,1]$时$t^k-t^{k+1}\geq0$，这些差的和为$t^n-t^m$，所以右侧绝对值不超过$\varepsilon t^m+\varepsilon(t^n-t^m)=\varepsilon t^n\leq\varepsilon$。
这一估计对所有$t\in[0,1]$同时成立，由一致柯西准则即得一致收敛。
\end{proof}
定理只要求端点处的数值级数收敛，允许它条件收敛；结论也只涉及收敛的那一端，不能推到发散的另一端。

\paragraph{和函数的连续性}
把这两个结论与上一节的连续性定理结合，就得到和函数的连续性。幂级数的各项是多项式，都连续。收敛区间内的每一点都落在某个内部闭子区间中，部分和在那里一致收敛，所以和函数在该点连续。

若$0<R<\infty$且某个端点包含在收敛域中，Abel第二定理给出从中心到该端点的一致收敛，和函数在该端点单侧连续（端点只能从区间内部趋近）。$R=0$时收敛域只有中心一点，单点定义域上的函数自动连续。
所以幂级数的和函数在实际收敛域内处处连续。

两个例子标出了结论的范围。几何级数在两个端点都发散，端点不在收敛域中，结论自然不涉及它们；事实上$x\to1^-$时和函数$1/(1-x)\to+\infty$，无法延伸到端点。
$\sum x^k/k$在$x=-1$处收敛，由Abel第二定理，它的和函数在$[-1,1)$上连续，在$-1$处右连续。第~\ref{sec:ru-elementary-taylor}节将证明这个和函数是$-\log(1-x)$，由$-1$处的连续性便得到交错调和级数的和为$\log2$。
